\documentclass{article}
\usepackage[utf8]{inputenc}
\usepackage{graphicx}
\usepackage{amsmath}
\usepackage{hyperref}
\usepackage{stackengine,amssymb,graphicx}
\newcommand\DDiamond{\ensurestackMath{%
  \stackengine{.5pt}{\Diamond}{\scalebox{.75}[1]{$-$}}{O}{c}{F}{F}{L}}}

\newcommand{\sol}{\textbf{solution}}

\title{CSE105/209A: Modern Algorithmic Toolbox\footnote{PI: Vaggos Chatziafratis}, PSET-2\\ }
\author{Ian Chi 2253175}
\date{}
\begin{document}
\maketitle

\section{}
\begin{enumerate}
  \item 
  \begin{equation}
    A = \begin{pmatrix}
    0 & 1 & 0 & 0 \\
    1 & 0 & 1 & 0 \\
    0 & 1 & 0 & 1 \\
    0 & 0 & 1 & 0
    \end{pmatrix}
  \end{equation}
  
  subtract \(\lambda\) from diagonal
  
  \begin{equation}
    A-\lambda I = \begin{pmatrix}
    -\lambda & 1 & 0 & 0 \\
    1 & -\lambda & 1 & 0 \\
    0 & 1 & -\lambda & 1 \\
    0 & 0 & 1 & -\lambda
    \end{pmatrix}
  \end{equation}
  
  the characteristic polynomial is 
  
  \begin{equation}
    \lambda(\lambda(\lambda^2-1)-\lambda) + (-1(\lambda^2-1)) = \lambda^4-3\lambda^2+1 = (\lambda^2 - \lambda - 1)(\lambda^2 + \lambda-1)
  \end{equation}
  which gives roots of 
  
  \begin{equation}
    \begin{split}
      \lambda_1 &= -\frac{1}{2}-\frac{\sqrt{5}}{2}\\
      \lambda_2 &= \frac{1}{2}-\frac{\sqrt{5}}{2}\\
      \lambda_3 &= -\frac{1}{2}+\frac{\sqrt{5}}{2}\\
      \lambda_4 &= \frac{1}{2}+\frac{\sqrt{5}}{2}
    \end{split}
  \end{equation}

  to find eigenvalues, find the null space after subtracting \(\lambda I\) from the diagonal of \(A\). 


  \begin{equation}
    (A-\lambda I)x=0 = \begin{pmatrix}
    -\lambda & 1 & 0 & 0 \\
    1 & -\lambda & 1 & 0 \\
    0 & 1 & -\lambda & 1 \\
    0 & 0 & 1 & -\lambda
    \end{pmatrix}
    x = 0
  \end{equation}
  
  \begin{equation}\begin{split}
  -\lambda x_1 + x_2 &= 0\\
  x_1 -\lambda x_2 + x_3 &= 0\\
  x_2 -\lambda x_3 + x_4 &= 0\\
  x_3 - \lambda x_4  &= 0\\
  \end{split}
  \end{equation}

  if we set \(x_1\) to 1, 

  \begin{equation}\begin{split}
      x_2 &= \lambda\\
      x_3 &= \lambda^2 - 1=(\lambda-1)(\lambda+1)\\
      x_4 &= \lambda - 1/\lambda
  \end{split}
  \end{equation}

  note

  \begin{equation}
    \left[\frac{1}{2}(\pm1 \pm \sqrt{5}) - 1\right] \left[\frac{1}{2}(\pm1 \pm \sqrt{5}) + 1\right] = \frac{1}{2}(\mp1 \mp \sqrt{5}) - 1
  \end{equation}
  \begin{equation}
    \left[\frac{1}{2}(\pm1 \mp \sqrt{5}) - 1\right] \left[\frac{1}{2}(\pm1 \mp \sqrt{5}) + 1\right] = \frac{1}{2}(\mp1 \pm \sqrt{5}) - 1
  \end{equation}

  \begin{equation}
    \left[\frac{1}{2}(1 \pm \sqrt{5}) - 1\right] - \frac{1}{\left[\frac{1}{2}(1 \pm \sqrt{5}) - 1\right]} = 1
  \end{equation}

  \begin{equation}
    \left[\frac{1}{2}(-1 \pm \sqrt{5}) - 1\right] - \frac{1}{\left[\frac{1}{2}(-1 \pm \sqrt{5}) - 1\right]} = -1
  \end{equation}

  for \(\lambda_1,\) 

  \begin{equation}
    \begin{split}
      v_1 &= \left( 1, -\frac{1}{2}-\frac{\sqrt{5}}{2}, \frac{1}{2}+\frac{\sqrt{5}}{2} , -1 \right)\\
      v_2 &= \left( 1, \frac{1}{2}-\frac{\sqrt{5}}{2}, -\frac{1}{2}+\frac{\sqrt{5}}{2} , 1 \right)\\
      v_3 &= \left( 1, -\frac{1}{2}+\frac{\sqrt{5}}{2}, \frac{1}{2}-\frac{\sqrt{5}}{2} , -1 \right)\\
      v_3 &= \left( 1, \frac{1}{2}+\frac{\sqrt{5}}{2}, -\frac{1}{2}-\frac{\sqrt{5}}{2} , 1 \right)\\
    \end{split}
  \end{equation}
  
  
  \newpage
  \item 
  \begin{equation}
    B = \begin{pmatrix}
    1 & -1 & 0 & 0 \\
    -1 & 2 & -1 & 0 \\
    0 & -1 & 2 & -1 \\
    0 & 0 & -1 & 1
    \end{pmatrix}
  \end{equation}

  \begin{equation}
    B-\lambda I = \begin{pmatrix}
    1-\lambda & -1 & 0 & 0 \\
    -1 & 2-\lambda & -1 & 0 \\
    0 & -1 & 2-\lambda & -1 \\
    0 & 0 & -1 & 1-\lambda
    \end{pmatrix}
  \end{equation}

  the characteristic polynomial of B is 

  \begin{equation}
    (1-\lambda)((2-\lambda)((2-\lambda)(1-\lambda) - 1) -(1-\lambda)) + ((-(2-\lambda)(1-\lambda))+1)
  \end{equation}
  which is equal to \(x(x-2)(x^2-4x+2)\) which gives 

  \begin{equation}
\begin{split}
    \lambda_0 &= 0\\
    \lambda_1  &= 2-\sqrt{2}\\
    \lambda_2 &= 2\\
    \lambda_3 &= 2+\sqrt{2}
\end{split}  \end{equation}

for eigenvectors, we know \(v_0=(1,1,1,1)\)

we can also guess \(v_2 = (1,-1,-1,1)\) and it verifies to be correct. 

the final two we can also guess, the diagonal is \(1\pm\sqrt{2}, \pm  \sqrt{2},\pm \sqrt{2}, 1\pm \sqrt{2}\)

\begin{equation}
  v_1 = (-1,1-\sqrt{2},-1+\sqrt{2},1)
\end{equation}

\begin{equation}
  v_3 = (-1,1+\sqrt{2},-1-\sqrt{2},1)
\end{equation}

\item 

  \begin{equation}
    C = \begin{pmatrix}
    0 & 1 & 1 & 1 \\
    1 & 0 & 1 & 1 \\
    1 & 1 & 0 & 1 \\
    1 & 1 & 1 & 0
    \end{pmatrix}
  \end{equation}

  one can guess the main eigenvector \((1,1,1,1)\) with eigenvalue 3. 

  
\end{enumerate}





\end{document}

